%---------------------------------------------------------------------------------
%                西南交通大学研究生学位论文：中文摘要
%---------------------------------------------------------------------------------

% 中文摘要输入区A
\cabstract
{
面向数字政府建设背景下持续积累的结构化政务数据，本文围绕非技术人员在政务场景中普遍面临的“查数难、分析难、协同难”问题，研究大语言模型驱动的数据查询与分析方法。针对政务术语与物理字段映射复杂、统计口径与计算逻辑隐含、查询与分析职责边界不清等特点，本文以“可靠查询、规范分析、系统落地”为总体目标，构建面向政务领域的智能问数与数据分析技术链路。

围绕上述目标，本文提出了基于逻辑增强与异构协同的政务数据查询生成方法，通过构建领域知识库、动态样例库与逻辑增强模式LA-Schema，将业务术语映射、逻辑计算公式和值域约束显式注入查询生成过程，并结合逻辑映射生成器、渐进分治生成器、执行反馈修复与候选判别模型，实现面向复杂政务业务语义的可靠查询生成。进一步地，本文提出了基于分析规划与查询增强的政务数据分析方法，引入分析规划模式AP-Schema，将开放式分析需求转化为结构化规划，通过查询增强的任务分解机制复用可靠取数能力，并结合轻量分析算子链与可视化规则生成规范化分析结果。在此基础上，本文设计并实现了一个面向政务场景的智能问数系统，将查询方法与分析方法组织为可交互、可追溯的统一服务。

实验结果表明，本文方法在自建政务数据集GovSQL上取得了91.47\%的执行准确率，较最强基线XiYan-SQL提升7.64个百分点；在56个政务分析任务上，流程成功率、任务完成率和结果正确率分别达到92.9\%、85.7\%和95.8\%，能够稳定生成图表、表格与文字洞察；所实现系统支持自然语言问数、分析结果生成与过程追溯，并在BIRD与Spider公共数据集上保持较强竞争力。

研究表明，本文形成了面向政务领域从自然语言需求理解、可靠取数到规范分析与系统实现的完整技术链路，能够为非技术人员开展政务问数与分析提供方法基础与系统支撑，也为大语言模型在垂直领域数据智能应用中的研究与落地提供了参考。
}
{政务数据，大语言模型，Text-to-SQL，智能问数，分析规划，数据分析}
% 中文关键词输入区，基于2024年新规范，关键词之间使用逗号“，”隔开，3~8个
